\chapter{Trabajo Realizado}
\label{cap:trabajo_realizado}
En el desarrollo de esta capitulo, desglocaremos las actividades a realizar, partiendo desde las reuniones con la Defensoría de Derechos Politécnicos, hasta llegar a la creación de un prototipo para el clasficador de competencia.

Las reuniones con la DDP fueron de gran importancia pues con solo leer los procesos ubicados en la documentación no bastaban, fue por esa razón que una junta con los dirigentes, permitio darle claridad a dudas que impedian darle continuidad al proyecto.

Despues de la primer junta, comenzamos con el desarrollo de vistas, acción que resulto contraproducente, pues si bien, las dudas inciales no existian, ahora nos encontrabamos con falta de acciones realizadas y como manejaban ciertas situaciones, así que regresamos a una segunda reunión, sin embargo, en esta segunda reunión contabamos con un diagrama BPMN, el cual se obtuvo por medio de sus documentos oficiales más el contexto que ya teniamos con anterioridad, al llegar a esa segunda reunión mencionaron que todo el proceso en papel no se seguia de manera manual, así que fuimos eliminando juntos con ellos cada parte innecesario, hasta obtener el siguiente diagrama BPMN:



Este diagrama representa el flujo de trabajo dentro de la defensoría, en donde se muestra como una queja pasa por una gran variedad de procesos, el flujo comienza en la recepción de la queja, una vez aceptada, se turna a un abogado disponible o al abogado que ya maneja este expediente, este abogado se encargara de definir la competencia o incompetencia, en caso de ser incompetente se le avisa al quejoso y se le informa cual es el lugar adecuado para su queja, si es competente pasa a subdefensoría a solicitar investigación, si la invetsigación no resulta satisfactoria se generan recordatorios y si se da respuesta a estos actos de investigación pasamos a un acuerdo de concilicación el cual se esperaria que fuera aprovado, si es aprovado se envia a la defensora con el fin de sellar y al ultimo a recepción para guardar en archivo.

Este proceso es meramente manual, pero seguían las dudas, así que una vez a la semana, nos vimos en la necesidad de agendar una cita por áreas separadas, comenzando con Repcionistas quienes nos explicaron a detalle las acciones que hacen y como podriamos mejorar su proceso. Despues de la reunión comenzamos con el desarrollo de los casos de uso y vistas para esta área.

Para la siguiente reunión nos dedicamos al área de primer contacto y ocurrio extamente lo mismo, explicaicón del proceso, creamos CU y vistas

De nuevo tuvimos una carta reunión en donde mostramos a los dos departamentos anteriores las vistas 